\subsection{分次环的有限性判别法}

在本节及以下诸节中，所考虑的所有分次（《代数》第 II 章 § 11）都假设是 Z 型的。若 A（相应地 M）是一个分次环（相应地分次模），则 \(A_{i}\) （相应地 \(M_{i}\) ）表示 A（相应地 M）中次数为 i 的齐次元组成的集合。

若 \(A_{i}=\{0\}\) （相应地 \(M_{i}=\{0\}\) ）对 i\textless0 成立，则为简便起见，称 A（相应地 M）为正次数的分次环（相应地分次模）。

命题 1. 设 \(\mathbf{A} = \bigoplus_{i\in \mathbb{Z}}\mathbf{A}_i\) 是一个正次数的分次交换环，\(\mathfrak{m}\) 是分次理想 \(\bigoplus_{i\geqslant 1}\mathbf{A}_i\) ，而 \((x_{\lambda})_{\lambda \in \mathbf{L}}\) 是 \(\mathbf{A}\) 的次数 \(\geqslant 1\) 的齐次元组成的一个族。下列条件等价：

\begin{enumerate}
\def\labelenumi{(\alph{enumi})}
\item
  \(\mathbf{A}\) 的由族 \((x_{\lambda})\) 生成的理想等于 \(\mathfrak{m}\) 。
\item
  族 \((x_{\lambda})\) 是 \(\mathbf{A}_0\) -代数 \(\mathbf{A}\) 的一个生成系。
\item
  对所有 \(i \geqslant 0\) ，\(\mathbf{A}_0\) -模 \(\mathbf{A}_i\) 由形如 \(\prod_{\lambda} x_{\lambda}^{n_{\lambda}}\) 且次数为 \(i\) （在 \(\mathbf{A}\) 中）的元素生成。
\end{enumerate}

显然条件 (b) 与 (c) 等价。若它们成立，则 m 的每个元素都形如 \(f((x_{\lambda}))\) ，其中 f 是 \(A_{0}[X_{\lambda}]_{\lambda \in L}\) 的一个无常数项的多项式；于是 \(m = \sum_{\lambda \in L} Ax_{\lambda}\) ，这就证明了 (c) 蕴含 (a)。反之，设条件 (a) 成立。设 \(A' = A_{0}[x_{\lambda}]_{\lambda \in L}\) 是 A 的子 \(A_{0}\) -代数，由族 \((x_{\lambda})\) 生成，我们来证明 \(A' = A\) 。为此，只需证明 \(A_{i} \subset A'\) 对所有 \(i \geqslant 0\) 成立。我们对 i 作归纳，性质对 i = 0 是显然的。设 \(y \in A_{i}\) ，其中 \(i \geqslant 1\) 。由于 \(y \in m\) ，故

存在一个有限支撑族 \((a_{\lambda})_{\lambda \in \mathbf{L}}\) ，其元素属于 \(\mathbf{A}\) ，使得 \(y = \sum_{\lambda} a_{\lambda} x_{\lambda}\) ，并且可以假设每个 \(a_{\lambda}\) 都是次数为 \(i - \deg(x_{\lambda})\) 的齐次元（必要时用该次数的齐次分量代替它）；由于 \(\deg(x_{\lambda}) > 0\) ，归纳假设表明 \(a_{\lambda} \in \mathbf{A}'\) 对所有 \(\lambda \in \mathbf{L}\) 成立，从而 \(y \in \mathbf{A}'\) 且 \(\mathbf{A}_t \subset \mathbf{A}'\) ，这就证明了 (a) 蕴含 (b)。

推论. 设 \(\mathbf{A} = \bigoplus_{i\in \mathbb{Z}}\mathbf{A}_i\) 是一个正次数的分次交换环，m 是分次理想 \(\bigoplus_{i\geqslant 1}\mathbf{A}_i\) 。

\begin{enumerate}
\def\labelenumi{(\roman{enumi})}
\tightlist
\item
  下列条件等价：
\end{enumerate}

\begin{enumerate}
\def\labelenumi{(\alph{enumi})}
\item
  理想 \(\mathfrak{m}\) 是一个有限生成 A-模。
\item
  环 A 是一个有限生成 \(A_{0}\) -模。
\end{enumerate}

\begin{enumerate}
\def\labelenumi{(\roman{enumi})}
\setcounter{enumi}{1}
\item
  设 (i) 中的条件成立，并设 \(\mathbf{M} = \bigoplus_{i\in \mathbf{Z}}\mathbf{M}_i\) 是一个有限生成分次 A-模。则对所有 \(i\in \mathbf{Z}\) ，\(\mathbf{M}_i\) 是一个有限生成 \(\mathbf{A}_0\) -模，并且存在 \(i_0\) ，使得 \(\mathbf{M}_i = \{0\}\) 对 \(i < i_0\) 成立。
\item
  若 A 的元素的一个族 \((y_{\mu})\) 是 A-模 m（相应地 \(A_{0}\) -模 A）的一个生成系，则由诸 \(y_{\mu}\) 的齐次分量组成的族也是如此；于是由命题 1 即得条件 (a) 与 (b) 的等价性。
\item
  我们可以假设 A（作为 \(\mathbf{A}_0\) -代数）由齐次元 \(a_i\) （ \(1 \leqslant i \leqslant r\) ，次数 \(\geqslant 1\) ）生成，且 M（作为 A-模）由齐次元 \(x_j\) （ \(1 \leqslant j \leqslant s\) ）生成；设 \(h_i = \deg(a_i)\) ，\(k_j = \deg(x_j)\) 。显然 \(\mathbf{M}_n\) 由系数在 \(\mathbf{A}_0\) 中的、形如 \(a_1^{\alpha_1} a_2^{\alpha_2} \ldots a_r^{\alpha_r} x_j\) 的元素的线性组合组成，其中诸 \(\alpha_i\) 是整数 \(\geqslant 0\) ，满足关系 \(k_j + \sum_{i=1} \alpha_i h_i = n\) ；对每个 \(n\) ，满足这些条件的族 \((\alpha_i)_{1 \leqslant i \leqslant r}\) 只有有限多个，因为 \(h_i \geqslant 1\) 对所有 \(i\) 成立；由此断言 \(\mathbf{M}_n\) 是一个有限生成 \(\mathbf{A}_0\) -模，并且显然 \(\mathbf{M}_n = \{0\}\) （当 \(n < \inf_{j}(k_j)\) 时）。
\end{enumerate}

